\documentclass[11pt]{article}
\usepackage[margin=1in]{geometry}
\usepackage{amsmath,amssymb}
\usepackage[hidelinks]{hyperref}

\begin{document}

\section*{Human Review: \(\Pi(7,15)=7/3\) via Equiangular Lines}

\noindent Original automatic proof: \texttt{solution\_packet.pdf}

\noindent Checked date: 15 June 2026

\noindent Status: Correct, but probably minor as a standalone observation.

\noindent Verdict: The proof appears correct. It gives a valid answer to Basso's printed question by combining the Koenig--Lewis--Lin equality criterion with a standard equiangular-line configuration in \(\mathbb R^7\). However, the mathematical novelty is likely limited if the underlying \(28\)-line construction is regarded as standard.

\section*{Human Comments}

The argument correctly identifies that the value \(7/3\) is obtained from the Koenig--Lewis--Lin bound at
\[
(n,d)=(7,15).
\]
Indeed,
\[
\frac7{15}
+
\sqrt{14\cdot\frac7{15}\cdot\frac8{15}}
=
\frac7{15}+\frac{28}{15}
=
\frac73.
\]
The quoted Koenig--Lewis--Lin theorem gives the upper bound, and its equality case is characterized by the existence of \(d\) equiangular lines in \(\mathbb R^n\). Thus it remains only to produce \(15\) equiangular lines in \(\mathbb R^7\).

The packet does this by using the standard \(28\)-line configuration in the hyperplane
\[
H=\left\{x\in\mathbb R^8:\sum_{i=1}^8 x_i=0\right\},
\]
where the lines are represented by
\[
v_{\{i,j\}}
=
e_i+e_j-\frac14(1,\ldots,1).
\]
For two-element sets \(A,B\subset\{1,\ldots,8\}\), one has
\[
\langle v_A,v_B\rangle=|A\cap B|-\frac12.
\]
Hence all these vectors have the same norm, and distinct lines have constant absolute inner product. Taking any \(15\) of these lines which span \(H\) gives the required equality case for \((n,d)=(7,15)\).

The selected \(15\) lines in the packet do span \(H\), since the seven vectors \(v_{\{1,j\}}\), \(2\leq j\leq 8\), already form a basis. Therefore the proof of
\[
\Pi(7,15)=\frac73
\]
looks correct.

\section*{Mathematical Value}

Although correct, this should probably be classified as a relatively minor observation rather than a substantial new result. The key geometric input is the classical \(28\)-line equiangular configuration in \(\mathbb R^7\), and the application to Basso's question is then immediate from the Koenig--Lewis--Lin equality criterion. In other words, the packet appears to notice a short consequence of known ingredients rather than introduce a new construction.

\section*{Recommended Follow-Up}

This example suggests a good next step for the automated pipeline. Rather than stopping at the isolated value
\[
\frac73,
\]
it would be useful to rerun or adapt the search to look for a more systematic pattern. In particular, one should try to find other pairs \((n,d)\) for which the Koenig--Lewis--Lin bound gives natural subsequent values, and then check whether known equiangular-line systems realize equality.

Ideally, the model should attempt to produce either:
\[
\text{a general form for such examples,}
\]
or
\[
\text{an infinite family of pairs }(n,d)\text{ giving further explicit values.}
\]
Such a generalization would be much more mathematically valuable than the single observation \(\Pi(7,15)=7/3\).

\section*{Review Outcome}

Accepted as correct, but with low novelty as a standalone contribution. Recommended classification: correct short observation/consequence of standard equiangular-line theory. Recommended next action: rerun the problem with instructions to search for higher values, additional pairs \((n,d)\), and possible infinite families.

\end{document}
